% Generated draft: requires mathematical and source review.
\documentclass[11pt]{article}
\usepackage[margin=25mm]{geometry}
\usepackage{fontspec}
\setmainfont{Latin Modern Roman}
\usepackage{amsmath,amssymb,mathtools,unicode-math}
\setmathfont{Latin Modern Math}
\usepackage{enumitem,needspace,xurl,hyperref}
\setlength{\parindent}{0pt}
\setlength{\parskip}{5pt}
\setlength{\emergencystretch}{3em}
\newcommand{\N}{\mathbb{N}}
\newcommand{\Z}{\mathbb{Z}}
\newcommand{\Q}{\mathbb{Q}}
\newcommand{\R}{\mathbb{R}}
\newcommand{\C}{\mathbb{C}}
\newcommand{\F}{\mathbb{F}}
\newcommand{\A}{\mathbb{A}}
\newcommand{\PP}{\mathbb{P}}
\newcommand{\E}{\mathbb{E}}
\newcommand{\eps}{\varepsilon}
\newcommand{\dd}{\,\mathrm{d}}
\begin{document}
\section*{UnsolvedMath ID: 800792}
\textbf{Source problem code:} catalogue entry
\subsection{Definitions and mathematical statement}
A holomorphic self-map germ fixing the origin is an equivalence class of holomorphic maps defined near \(0\in\C^3\), written
\[
f:(\C^3,0)\longrightarrow(\C^3,0),\qquad f(0)=0,
\]
where two representatives agree on a neighborhood of the origin. Such a germ is dominant if \(\det Df\) is not identically zero, and it is non-invertible if it has no holomorphic inverse germ; equivalently, \(\det Df(0)=0\).

\textbf{Source classification problem.} Classify the three-dimensional non-invertible dominant rigid germs.

The supplied excerpt does not define \emph{rigid}, specify the equivalence relation for classification, or reproduce the hypotheses of the preceding results that it proposes to extend to dimension three. Consequently, the following gives a conditional interpretation, requiring confirmation against the source context.

For a dominant germ \(f\), define its critical set and generalized critical set by
\[
\operatorname{Crit}(f)=\{z:\det Df(z)=0\},\qquad
\operatorname{Crit}_{\infty}(f)=\bigcup_{n\geq 0}f^{-n}\bigl(\operatorname{Crit}(f)\bigr),
\]
with inverse images and unions understood at the level of germs at the origin. One usual rigidity condition is that this union is a reduced normal-crossings divisor: in some local holomorphic coordinates, its set germ is
\[
\{z_1\cdots z_r=0\}\qquad\text{for some }r\in\{1,2,3\}.
\]
If this is the intended rigidity condition and classification means classification under local holomorphic conjugacy, the task is to classify all dominant, non-invertible germs satisfying this condition under
\[
f\sim g
\quad\Longleftrightarrow\quad
\text{there exists a biholomorphic germ }h:(\C^3,0)\to(\C^3,0)
\text{ with }h\circ f=g\circ h.
\]
A classification would give representatives and a necessary and sufficient criterion for two representatives to be equivalent. The supplied excerpt does not establish this choice of rigidity or equivalence, nor whether additional attracting or contracting hypotheses are inherited from the preceding discussion; these qualifications remain unresolved.
\subsection{Short English statement}
Abate's Problem OP11 asks for a classification of three-dimensional non-invertible dominant rigid holomorphic germs. The supplied excerpt leaves the definition of rigidity, the classification equivalence, and any inherited dynamical hypotheses unspecified. The expanded formulation above is therefore conditional. The registry labels the problem open; this is supplied source information, not an independently verified assessment.
\subsection{Sources}
\begin{itemize}
\item Local UnsolvedMath registry, supplied complete record: permanent ID \(800792\), source problem code \texttt{catalogue entry}, title ``Local Discrete Holomorphic Dynamics --- Open Problem OP11.'' The record attributes the problem to Abate, with proposed year 2011, and supplies the status label ``open.''
\item Abate, \emph{Open Problems in Local Discrete Holomorphic Dynamics} (2011), Problem OP11. Supplied source URL: \url{https://arxiv.org/abs/1107.4876}. The supplied provenance records access on 2026-10-16 and extraction method \texttt{arxiv\_latex}; these provenance claims have not been independently verified here.
\item The supplied introductory sentence presents this problem as an extension of preceding results to dimension three and as one of two additional open problems. Those preceding results and their hypotheses are not included in the supplied record.
\end{itemize}
\end{document}
